\documentclass[aspectratio=169,10pt]{beamer}

% Компилировать XeLaTeX или LuaLaTeX (fontspec + polyglossia)
\usepackage{fontspec}
\usepackage{polyglossia}
\setmainlanguage{russian}
\setotherlanguage{english}
\setmainfont{DejaVu Sans}[ItalicFont={DejaVu Sans},ItalicFeatures={FakeSlant=0.22},BoldItalicFont={DejaVu Sans},BoldItalicFeatures={FakeSlant=0.22}]
\setsansfont{DejaVu Sans}[ItalicFont={DejaVu Sans},ItalicFeatures={FakeSlant=0.22},BoldItalicFont={DejaVu Sans},BoldItalicFeatures={FakeSlant=0.22}]
\setmonofont{DejaVu Sans Mono}[Scale=0.88]
\newfontfamily\cyrillicfont{DejaVu Sans}[ItalicFont={DejaVu Sans},ItalicFeatures={FakeSlant=0.22},BoldItalicFont={DejaVu Sans},BoldItalicFeatures={FakeSlant=0.22}]
\newfontfamily\cyrillicfontsf{DejaVu Sans}[ItalicFont={DejaVu Sans},ItalicFeatures={FakeSlant=0.22},BoldItalicFont={DejaVu Sans},BoldItalicFeatures={FakeSlant=0.22}]
\newfontfamily\cyrillicfonttt{DejaVu Sans Mono}[Scale=0.88]

\usepackage{graphicx}
\usepackage{tikz}
\usetikzlibrary{arrows.meta,shapes.multipart,calc,positioning,fit}

\definecolor{main}{HTML}{1F4E79}
\definecolor{accent}{HTML}{B5472B}
\setbeamercolor{structure}{fg=main}
\setbeamercolor{frametitle}{fg=white,bg=main}
\setbeamercolor{block title}{fg=white,bg=main}
\setbeamercolor{block body}{bg=main!7}
\setbeamercolor{block title alerted}{fg=white,bg=accent}
\setbeamercolor{block body alerted}{bg=accent!8}
\setbeamerfont{frametitle}{size=\large,series=\bfseries}
\setbeamerfont{block title}{size=\small,series=\bfseries}
\setbeamerfont{block body}{size=\small}
\setbeamerfont{block body alerted}{size=\small}
\setbeamerfont{itemize/enumerate body}{size=\small}
\setbeamerfont{itemize/enumerate subbody}{size=\footnotesize}
\setbeamertemplate{navigation symbols}{}
\setbeamertemplate{items}[circle]
\setbeamersize{text margin left=0.8cm,text margin right=0.8cm}
\setbeamertemplate{frametitle}{%
  \nointerlineskip
  \begin{beamercolorbox}[wd=\paperwidth,ht=1.05cm,dp=0.35cm,leftskip=0.8cm,rightskip=0.8cm]{frametitle}
    \usebeamerfont{frametitle}\insertframetitle
  \end{beamercolorbox}%
  \vskip1.5mm
}
\setbeamertemplate{footline}{%
  \hbox to\paperwidth{\hspace{0.8cm}\color{black!55}\scriptsize
    Структурные паттерны проектирования\hfill
    \insertframenumber\,/\,\inserttotalframenumber\hspace{0.8cm}}%
  \vskip3mm
}

\tikzset{
  uml/.style={draw=black!70, fill=main!6, rectangle split,
    rectangle split draw splits, font=\footnotesize,
    text width=3.2cm, inner sep=3pt},
  cls3/.style={uml, rectangle split parts=3,
    rectangle split part align={center,left,left}},
  cls2/.style={uml, rectangle split parts=2,
    rectangle split part align={center,left}},
  cls1/.style={uml, rectangle split parts=1, align=center},
  % уменьшенный шрифт: применяется ПОСЛЕ cls*, без глобальных побочных эффектов
  scls1/.style={cls1, font=\scriptsize},
  scls2/.style={cls2, font=\scriptsize},
  scls3/.style={cls3, font=\scriptsize},
  % обобщение: сплошная линия, незакрашенный треугольник
  inherit/.style={-{Triangle[open,length=9pt,width=8pt]}, thick},
  % реализация: пунктир, незакрашенный треугольник
  realize/.style={-{Triangle[open,length=9pt,width=8pt]}, thick, dashed},
  % навигируемая ассоциация: сплошная линия, открытая «галочка» (V)
  assoc/.style={thick, -{Straight Barb[length=8pt,width=6pt]}},
  % зависимость: пунктир, открытая «галочка» (V)
  depend/.style={-{Straight Barb[length=9pt,width=7pt]}, thick, dashed},
  % композиция/агрегация: ромб у целого
  compo/.style={{Diamond[length=10pt,width=7pt]}-, thick},
  ndash/.style={dashed, black!60},
  lbl/.style={font=\scriptsize, inner sep=1.5pt, fill=white, text=black!80, align=center},
  note/.style={draw=black!50, fill=yellow!22, font=\scriptsize\ttfamily,
    align=left, inner sep=3pt},
}

\title{Структурные паттерны проектирования}
\date{\today}

\begin{document}

\begin{frame}[plain]
  \vfill
  \centering
  {\color{main}\fontsize{26}{31}\selectfont\bfseries Структурные паттерны\\проектирования\par}
  \vspace{5mm}
  {\color{accent}\rule{6cm}{1.2pt}\par}
  \vspace{5mm}
  {\large UML-диаграммы\par}
  \vspace{7mm}
  {\small Adapter \,·\, Bridge \,·\, Composite \,·\, Decorator \,·\, Facade, \\[1mm]
  Proxy\par}
  \vfill
\end{frame}

\begin{frame}{Содержание}
  \begin{columns}[T]
    \begin{column}{0.46\textwidth}
      \begin{enumerate}\setlength\itemsep{3pt}
        \item Адаптер (Adapter)
        \item Мост (Bridge)
        \item Компоновщик (Composite)
        \item Декоратор (Decorator)
      \end{enumerate}
    \end{column}
    \begin{column}{0.46\textwidth}
      \begin{enumerate}\setcounter{enumi}{4}\setlength\itemsep{3pt}
        \item Фасад (Facade)
        \item Заместитель (Proxy)
      \end{enumerate}
    \end{column}
  \end{columns}
\end{frame}

% --------------------------------------------------------------- Слайд 3
\begin{frame}{Adapter: UML-диаграмма}
  \centering
  \begin{tikzpicture}
    \node[cls1, text width=1.4cm] (client) at (0.0,3.4) {\textbf{Client}};
    \node[cls2, text width=3.0cm] (t) at (5.8,3.4) {«interface»\\\textbf{Target}
      \nodepart{two}+ request()};
    \node[cls2, text width=3.2cm] (a) at (5.8,0.4) {\textbf{Adapter}
      \nodepart{two}− adaptee : Adaptee\\+ request()};
    \node[cls2, text width=3.0cm] (ad) at (11.2,0.4) {\textbf{Adaptee}
      \nodepart{two}+ specificRequest()};
    \node[note] (n) at (1.5,0.4) {adaptee\\\quad .specificRequest();};
    \draw[depend] (client.east) -- (t.west)
      node[lbl, midway, above=1mm] {request()};
    \draw[realize] (a.north) -- (t.south);
    \draw[assoc] (a.east) -- (ad.west)
      node[lbl, midway, above=1mm] {adaptee};
    \draw[ndash] (n.east) -- (a.west);
  \end{tikzpicture}

  \vspace{3mm}
  {\small Клиент работает с интерфейсом Target. Adapter реализует его и переводит вызовы
   в вызовы несовместимого интерфейса класса Adaptee.}
\end{frame}

% --------------------------------------------------------------- Слайд 4
\begin{frame}{Bridge: UML-диаграмма}
  \centering
  \begin{tikzpicture}
    \node[cls1, text width=1.4cm] (client) at (-0.3,3.4) {\textbf{Client}};
    \node[cls3, text width=3.4cm] (ab) at (3.8,3.4) {\textbf{Abstraction}
      \nodepart{two}− impl : Implementor
      \nodepart{three}+ operation()};
    \node[cls1, text width=3.4cm] (ra) at (3.0,0.4) {\textbf{RefinedAbstraction}};
    \node[cls2, text width=3.0cm] (im) at (9.2,3.4) {«interface»\\\textbf{Implementor}
      \nodepart{two}+ operationImpl()};
    \node[scls1, text width=3.3cm] (ca) at (7.5,0.4) {\textbf{ConcreteImplementorA}};
    \node[scls1, text width=3.3cm] (cb) at (11.0,0.4) {\textbf{ConcreteImplementorB}};
    \draw[depend] (client.east) -- (ab.west);
    \draw[inherit] (ra.north) -- (ab.south);
    \draw[assoc] (ab.east) -- (im.west)
      node[lbl, midway, above=1mm] {impl};
    \draw[realize] ([xshift=-0.3cm]ca.north) -- ([xshift=-0.55cm]im.south);
    \draw[realize] ([xshift=0.3cm]cb.north) -- ([xshift=0.55cm]im.south);
  \end{tikzpicture}

  \vspace{3mm}
  {\small Абстракция и реализация образуют две независимые иерархии; связь между ними~--- ссылка
   \texttt{impl}, поэтому каждую иерархию можно развивать отдельно.}
\end{frame}

% --------------------------------------------------------------- Слайд 5
\begin{frame}{Composite: UML-диаграмма}
  \centering
  \begin{tikzpicture}
    \node[cls1, text width=1.4cm] (client) at (0.0,3.4) {\textbf{Client}};
    \node[cls2, text width=4.4cm] (c) at (6.4,3.4) {«interface»\\\textbf{Component}
      \nodepart{two}+ operation()\\+ add(c : Component)\\+ remove(c : Component)};
    \node[cls2, text width=2.4cm] (l) at (1.8,0.2) {\textbf{Leaf}
      \nodepart{two}+ operation()};
    \node[cls3, text width=4.0cm] (cp) at (8.6,0.0) {\textbf{Composite}
      \nodepart{two}− children : List<Component>
      \nodepart{three}+ operation()\\+ add(c : Component)\\+ remove(c : Component)};
    \draw[depend] (client.east) -- (c.west);
    \draw[realize] (l.north) -- ([xshift=-1.8cm]c.south);
    \draw[realize] ([xshift=-1.5cm]cp.north) -- ([xshift=0.5cm]c.south);
    \draw[compo] ([xshift=0.5cm]cp.north) -- ([xshift=2.2cm]c.south)
      node[lbl, pos=0.6, right=1mm] {children};
    \node[lbl, anchor=south west] at ([xshift=0.5cm,yshift=0.02cm]cp.north) {0..*};
  \end{tikzpicture}

  \vspace{1mm}
  {\small Лист и составной объект реализуют единый интерфейс Component, поэтому клиент
   одинаково работает с одиночными объектами и с целыми деревьями.}
\end{frame}

% --------------------------------------------------------------- Слайд 6
\begin{frame}{Decorator: UML-диаграмма}
  \centering
  \begin{tikzpicture}
    \node[cls1, text width=3.5cm] (cc) at (1.7,3.4) {\textbf{ConcreteComponent}};
    \node[cls2, text width=3.0cm] (c) at (6.0,3.4) {«interface»\\\textbf{Component}
      \nodepart{two}+ operation()};
    \node[cls1, text width=1.4cm] (client) at (11.2,3.4) {\textbf{Client}};
    \node[cls3, text width=3.4cm] (d) at (6.0,0.3) {\textbf{Decorator}
      \nodepart{two}− component : Component
      \nodepart{three}+ operation()};
    \node[cls1, text width=3.5cm] (da) at (1.7,0.3) {\textbf{ConcreteDecoratorA}};
    \node[cls1, text width=3.5cm] (db) at (10.7,0.3) {\textbf{ConcreteDecoratorB}};
    \draw[realize] (cc.east) -- (c.west);
    \draw[depend] (client.west) -- (c.east);
    \draw[realize] ([xshift=-0.7cm]d.north) -- ([xshift=-0.7cm]c.south);
    \draw[compo] ([xshift=0.9cm]d.north) -- ([xshift=0.9cm]c.south);
    \node[lbl, anchor=west] at ([xshift=0.95cm,yshift=1.1cm]d.north) {component};
    \draw[inherit] (da.east) -- (d.west);
    \draw[inherit] (db.west) -- (d.east);
  \end{tikzpicture}

  \vspace{3mm}
  {\small Декоратор реализует тот же интерфейс, что и оборачиваемый объект, и хранит ссылку на него;
   дополнительное поведение добавляется до или после делегирования вызова.}
\end{frame}

% --------------------------------------------------------------- Слайд 7
\begin{frame}{Facade: UML-диаграмма}
  \centering
  \begin{tikzpicture}
    \node[cls1, text width=1.4cm] (client) at (0.0,1.7) {\textbf{Client}};
    \node[cls2, text width=3.4cm] (f) at (4.6,1.7) {\textbf{Facade}
      \nodepart{two}+ operation()};
    \node[cls1, text width=2.2cm] (a) at (10.0,3.0) {\textbf{SubsystemA}};
    \node[cls1, text width=2.2cm] (b) at (10.0,1.7) {\textbf{SubsystemB}};
    \node[cls1, text width=2.2cm] (c) at (10.0,0.4) {\textbf{SubsystemC}};
    \node[draw=black!50, dashed, fit=(a)(b)(c), inner sep=9pt,
      label={[font=\scriptsize,text=black!70]below:Подсистема}] (sub) {};
    \draw[depend] (client.east) -- (f.west)
      node[lbl, midway, above=1mm] {operation()};
    \draw[depend] (f.east) -- (a.west);
    \draw[depend] (f.east) -- (b.west);
    \draw[depend] (f.east) -- (c.west);
  \end{tikzpicture}

  \vspace{3mm}
  {\small Facade предоставляет единый упрощённый интерфейс к набору классов подсистемы
   и скрывает от клиента их взаимодействие.}
\end{frame}

% % --------------------------------------------------------------- Слайд 8
% \begin{frame}{Front Controller (единая точка входа): UML-диаграмма}
%   \centering
%   \begin{tikzpicture}
%     \node[cls1, text width=1.4cm] (client) at (-0.3,2.0) {\textbf{Client}};
%     \node[cls3, text width=3.8cm] (fc) at (4.0,2.0) {\textbf{FrontController}
%       \nodepart{two}− commands : Map<String, Command>
%       \nodepart{three}+ handleRequest(r : Request)\\− authenticate(r : Request)};
%     \node[cls2, text width=2.8cm] (cmd) at (9.6,3.2) {«interface»\\\textbf{Command}
%       \nodepart{two}+ execute(r : Request)};
%     \node[scls1, text width=2.9cm] (c1) at (8.0,0.4) {\textbf{ConcreteCommandA}};
%     \node[scls1, text width=2.9cm] (c2) at (11.4,0.4) {\textbf{ConcreteCommandB}};
%     \draw[depend] (client.east) -- (fc.west)
%       node[lbl, midway, above=1mm] {request};
%     \draw[assoc] (fc.north east) -- (cmd.west)
%       node[lbl, midway, above=1mm] {execute()};
%     \draw[realize] ([xshift=-0.3cm]c1.north) -- ([xshift=-0.65cm]cmd.south);
%     \draw[realize] ([xshift=0.3cm]c2.north) -- ([xshift=0.65cm]cmd.south);
%   \end{tikzpicture}

%   \vspace{3mm}
%   {\small Все запросы поступают в один обработчик: он выполняет общую обработку (аутентификация,
%    журналирование), выбирает по запросу нужную команду и передаёт ей управление.}
% \end{frame}

% --------------------------------------------------------------- Слайд 10
\begin{frame}{Proxy: UML-диаграмма}
  \centering
  \begin{tikzpicture}
    \node[cls1, text width=1.4cm] (client) at (0.0,3.4) {\textbf{Client}};
    \node[cls2, text width=3.0cm] (s) at (6.2,3.4) {«interface»\\\textbf{Subject}
      \nodepart{two}+ request()};
    \node[cls2, text width=3.4cm] (p) at (3.0,0.3) {\textbf{Proxy}
      \nodepart{two}− realSubject : RealSubject\\+ request()};
    \node[cls2, text width=3.0cm] (rs) at (10.4,0.3) {\textbf{RealSubject}
      \nodepart{two}+ request()};
    \node[note] (n) at (6.7,-0.8) {if (checkAccess())\\\quad realSubject.request();};
    \draw[depend] (client.east) -- (s.west)
      node[lbl, midway, above=1mm] {request()};
    \draw[realize] ([xshift=0.3cm]p.north) -- ([xshift=-0.6cm]s.south);
    \draw[realize] ([xshift=-0.3cm]rs.north) -- ([xshift=0.6cm]s.south);
    \draw[assoc] (p.east) -- (rs.west)
      node[lbl, pos=0.5, above=1mm] {realSubject};
    \draw[ndash] (n.north) -- (6.7,0.3);
  \end{tikzpicture}

  \vspace{1mm}
  {\small Proxy имеет тот же интерфейс, что и RealSubject, и контролирует доступ к нему
   (защита, ленивое создание, кэширование, журналирование).}
\end{frame}

\end{document}